\documentclass[11pt,a4paper]{article}
\usepackage[T1]{fontenc}
\usepackage[utf8]{inputenc}
\usepackage{lmodern,microtype}
\usepackage[a4paper,margin=27mm,headheight=14pt]{geometry}
\usepackage{amsmath,amssymb,amsthm,mathtools}
\usepackage{booktabs,array,longtable}
\usepackage{xcolor}
\usepackage{enumitem}
\usepackage{fancyhdr}
\usepackage{xurl}
\usepackage[colorlinks=true,linkcolor=blue!40!black,citecolor=blue!40!black,urlcolor=blue!40!black]{hyperref}
\hypersetup{pdftitle={Moment Determinacy and Nonlinear Transseries},pdfauthor={Research report prepared for Vladimir Reshetnikov},pdfsubject={Positive Stieltjes realizations, a sharp Gevrey threshold, exact flat defects, and nonlinear inverse folds}}
\pagestyle{fancy}
\fancyhf{}
\fancyhead[L]{\small Moment Determinacy and Nonlinear Transseries}
\fancyhead[R]{\small Research report}
\fancyfoot[C]{\thepage}
\setlength{\emergencystretch}{2em}
\setlength{\parskip}{3pt}
\setcounter{tocdepth}{2}
\numberwithin{equation}{section}
\newtheorem{theorem}{Theorem}[section]
\newtheorem{proposition}[theorem]{Proposition}
\newtheorem{lemma}[theorem]{Lemma}
\newtheorem{corollary}[theorem]{Corollary}
\theoremstyle{definition}
\newtheorem{definition}[theorem]{Definition}
\newtheorem{example}[theorem]{Example}
\theoremstyle{remark}
\newtheorem{remark}[theorem]{Remark}
% ed. (2026-09-29): unnumbered environment for editorial notes added on filing (no numbering changes).
\newtheorem*{ednoteinner}{Editorial note (ProveIt, 2026-09-29)}
\newenvironment{ednote}{\begin{ednoteinner}\sloppy\Urlmuskip=0mu plus 1mu\relax}{\end{ednoteinner}}
\newcommand{\R}{\mathbb R}
\newcommand{\C}{\mathbb C}
\newcommand{\N}{\mathbb N}
\newcommand{\Z}{\mathbb Z}
\newcommand{\dd}{\,\mathrm d}
\newcommand{\ii}{\mathrm i}
\newcommand{\ee}{\mathrm e}
\newcommand{\coef}[2]{[#1^{#2}]}
\newcommand{\supp}{\operatorname{supp}}
\newcommand{\Res}{\operatorname{Res}}
\newcommand{\dist}{\operatorname{dist}}
\newcommand{\Id}{\operatorname{id}}
\newcommand{\Gev}{\operatorname{Gev}}
\newcommand{\cayley}[1]{\frac{#1^{#1-1}}{#1!}}
\title{\textbf{Moment Determinacy and\newline Nonlinear Transseries}\\[7pt]
\large A sharp Gevrey threshold, exact flat defects,\\
convergent ambiguity sectors, and Lambert--\(W\) folds}
\author{Research report prepared for Vladimir Reshetnikov\\
\small In connection with the ProveIt research program}
\date{29 September 2026}

\begin{document}
\maketitle
\thispagestyle{empty}
\begin{abstract}
A complete formal expansion need not identify the function that is to be
inverted, even when the function is a Stieltjes transform of a positive
probability measure. We study the nonlinear realization problem
\(T(y)=yS(T(y))\), where
\(S(t)=\int_0^\infty(1+t\lambda)^{-1}\dd\mu(\lambda)\).
First, inverse identifiability is shown to be exactly equivalent to
Stieltjes moment determinacy. A cancellation-free Lagrange formula
preserves the exact factorial growth order of the moment sequence.
Consequently, formal Gevrey order at most two forces uniqueness in this
positive-measure class, whereas every order greater than two admits
explicit continua of distinct realizations.

Two complementary families make the hidden data explicit. Generalized
gamma measures give an exact defect
\(Ct^{-\beta}\exp(-At^{-\alpha})\); bilateral geometric measures give a
reciprocal theta function with a nonconstant log-periodic amplitude.
For the latter, the ambiguity is comparable to the least moment-remainder
bound. We transport both defects through nonlinear inversion, calculate
the first relative correction, and construct convergent expansions in
successive powers of the flat defect with explicit Cauchy bounds.
After a defect-dependent rescaling, the deformation map converges to
\(v\ee^{-v}\). This proves an asymptotically exact radius for the
parameter expansion, Cayley-number limits for its coefficients, and a
nearby square-root fold, including its first gamma and log-periodic
corrections.

The moment-indeterminacy constructions, Hardy-type determinacy argument,
Jacobi identity, and Lagrange inversion are classical foundations. The
paper develops their quantitative nonlinear synthesis, without claiming
exhaustive historical priority or a solution of the repository's general
Hahn-series realization problem. All proofs are conventional mathematical
proofs, not Lean formalizations. Reproducible finite checks and a precise
research agenda accompany the article.
\end{abstract}
\vfill
\noindent\textbf{Scope.} Positive-real asymptotic germs as \(y\downarrow0\);
formal coefficient growth; explicit analytic realizations and inverse
parameter singularities. No general summation rule is assumed.
\clearpage
\tableofcontents
\clearpage

\section{The realization question and the contribution}
\label{sec:context}

\subsection{A focused question arising from the repository}
The ProveIt transseries corpus distinguishes formal algebra, asymptotic
scales, flat functions, and analytic inversion certificates
\cite{repo-main,repo-series}. These distinctions suggest a precise
realization question:
\begin{quote}
If every coefficient of an inverse expansion is known, and the
underlying function is required to come from a positive measure, is the
actual inverse uniquely determined? If not, what is the size and analytic
structure of the missing information?
\end{quote}
This question is narrower than the construction of a general field of
transseries and broader than computing another list of inverse
coefficients. It admits a complete answer at the level of identifiability,
a sharp coefficient-growth threshold, and explicit models of the
remaining ambiguity.

The inspected repository companion on action accumulation already
exhibits flat oscillations and studies their transport under inversion
\cite{repo-action}. Its editorial notes explicitly distinguish those
results from the canonical volume's general realization question,
identified there by the label
\texttt{plt:rmk:ext-open-realization}. We maintain that distinction.
Here the support is that of a \emph{positive Stieltjes moment problem};
we do not construct a realization of every left-finite Hahn series.
The positive measure is the additional structure, and moment
determinacy is exactly the extra condition that restores identifiability.

% ed. (2026-09-29): repository results on the determinate side, and two notational links, not cited in the delivered text (dossier of batch 51).
\begin{ednote}
Two filed repository results lie on the determinate side of
Theorem~\ref{thm:identifiability}. The canonical volume \cite{ed:tai}
(its Theorem \texttt{p8:thm:bell}) writes the tail
\(C(x)=\ee^{-1}\sum_{j\ge0}1/(j!\,(x+j+1))\) as a Stieltjes transform
whose moments are the Bell numbers, with a signed remainder. In the
present notation \(C(x)=x^{-1}S_\mu(1/x)\) for the probability measure
\(\mu=\sum_{j\ge0}\delta_{j+1}/(\ee\,j!)\), with moments
\(m_n=\mathfrak B_{n+1}\le(n+1)!\), of Gevrey order at most one; so
Lemma~\ref{lem:hardy} makes that representation the only positive one,
and by Theorem~\ref{thm:identifiability} the corresponding inverse has
exactly one positive Stieltjes realization. The certified \(q\to1\)
package \cite{ed:ciq} (its \texttt{thm:stieltjes} and \texttt{thm:pade})
writes a finite-size correction as \(\int\dd\nu_\tau(u)/(1+su)\) with a
positive measure whose moments are bounded by \(CA^{2j}(2j)!\), that is,
of factorial-square growth, the boundary case \(s=2\) of
Theorem~\ref{thm:gevrey}; it proves determinacy by the argument of
Lemma~\ref{lem:hardy} and adds two-sided Pad\'e bounds. The
inverse-digamma package \cite{ed:ids} has a Stieltjes-type representation
only modulo a convergent Laurent function, outside the class used here.
The word \emph{realization} in this article means a positive Stieltjes
representation of the inverse germ at \(y\downarrow0\), not the
Hardy-field sense of the canonical volume's
Definition~\texttt{plt:def:cmp-realization}. The limit
\(v_0=-W_0(-\kappa)\) of \eqref{eq:lambert-limit} is that volume's Cayley
tree function (\texttt{p1:def:cayley}). This note is editorial.
\end{ednote}

\subsection{The main results in one place}
Let \(m_n=\int\lambda^n\dd\mu\), \(m_0=1\), and write the formal inverse as
\begin{equation}
 \widehat T(y)=\sum_{n\ge1}b_ny^n,
 \qquad c_n=(-1)^n b_{n+1}.
\end{equation}
The following statements will be proved in order.

\paragraph{An exact identifiability criterion.}
The actual germs \(T_\mu\) and the representing measures \(\mu\) are in
one-to-one correspondence. All measures with the same moments produce
the same formal inverse. Thus that formal inverse identifies a unique
positive Stieltjes realization if and only if its moment problem is
determinate (Theorem~\ref{thm:identifiability}).

\paragraph{A sharp factorial-growth threshold.}
The coefficients satisfy the noncancelling identity
\begin{equation}
 c_n=\frac1{n+1}
 \sum_{k_1+\cdots+k_{n+1}=n}m_{k_1}\cdots m_{k_{n+1}}.
 \label{eq:intro-composition}
\end{equation}
In particular, \(m_n\le c_n\), and a bound
\(m_n\le CA^n(n!)^s\) implies
\(c_n\le(4AC)^n(n!)^s\). Formal Gevrey order \(s\le2\) therefore forces
uniqueness. For every \(s>2\), a generalized gamma family gives distinct
positive realizations of exact formal Gevrey order \(s\)
(Theorems~\ref{thm:gevrey} and~\ref{thm:gamma}).

\paragraph{Explicit missing data.}
The generalized gamma family has a closed exponential defect, while the
bilateral geometric family has a reciprocal theta defect. The latter
has a positive log-periodic amplitude and is comparable to the
optimally chosen moment-remainder bound
(Theorems~\ref{thm:gamma-defect},~\ref{thm:theta-defect},
and~\ref{thm:least}). Even an exact first-order \(q\)-difference equation,
added to positivity and the complete formal expansion, need not select a
unique transform (Proposition~\ref{prop:shifted}).

\paragraph{Nonlinear sectors and their radius.}
For a physical baseline \(t_0=T_{a_0}(y)\), put \(D_0=D(t_0)\) and let
\(p\to\infty\) be the logarithmic slope scale of the defect. The expansion
of \(T_{a_0+\delta}(y)\) in \(\delta\) converges, with explicit coefficient
and remainder bounds. Its radius \(R_y\) obeys
\begin{equation}
 R_y\sim\frac{t_0}{\ee\,p\,yD_0}.
 \label{eq:intro-radius}
\end{equation}
At fixed sector number \(k\), the normalized coefficient tends to
\(k^{k-1}/k!\). The critical deformation is a Lambert--\(W\) fold, and
its first correction retains the theta phase
(Theorems~\ref{thm:sectors},~\ref{thm:radius}, and~\ref{thm:fold-correction}).

\subsection{What is classical, and what is claimed here}
Moment-indeterminate generalized gamma and Stieltjes--Wigert measures
are classical; useful primary accounts include
\cite{lin-stoyanov,chihara,christiansen,christiansen-koelink}.
The order-two determinacy boundary is a Hardy/Carleman phenomenon,
not a newly discovered boundary in moment theory. Likewise, Jacobi's
triple product, Lagrange inversion, and the local square-root normal form
are classical tools. They are either stated with attribution or derived
in the form used below.

The contribution developed in this report is the organized nonlinear
realization theorem, its cancellation-free coefficient control, the
exact defect transport with relative corrections, the quantitative
convergent sector expansion, and the defect-dependent parameter-radius
and fold analysis. These are proposed research contributions with proofs;
we do not assert that no equivalent formulation has appeared elsewhere.
The source comparison was targeted, not a theorem-by-theorem audit of the
large repository corpus. No result in this report is claimed to have been
checked by Lean.

It is also important not to confuse the present nonuniqueness with a
failure of a sectorial implicit-function theorem. A prescribed
Borel-summable class contains additional analytic information; for
example, Nikolaev's theorem concerns a selected summable asymptotic
setting \cite{nikolaev}. Our input data are positivity and a full formal
jet, without such a selection rule. The various \(q\)-summation procedures
also have their own hypotheses \cite{divizio-zhang}; none is silently
identified with the families below.

\section{Positive Stieltjes transforms and their inverse germs}
\label{sec:framework}

\begin{definition}
Let \(\mathcal P_\infty\) be the probability measures on \([0,\infty)\)
with finite moments of every nonnegative integer order. For
\(\mu\in\mathcal P_\infty\), define
\begin{equation}
 S_\mu(t)=\int_0^\infty\frac{\dd\mu(\lambda)}{1+t\lambda},
 \qquad F_\mu(t)=\frac{t}{S_\mu(t)}
 \quad(t>0).
 \label{eq:stieltjes}
\end{equation}
The \emph{positive Stieltjes inverse} is \(T_\mu=F_\mu^{-1}\).
Thus its defining equation is
\begin{equation}
 T_\mu(y)=yS_\mu(T_\mu(y)).
 \label{eq:implicit}
\end{equation}
A flat function at zero is one that is \(O(t^N)\) for every \(N\).
\end{definition}

\begin{lemma}[Regularity, remainder, and positivity]
\label{lem:regularity}
The function \(S_\mu\) is holomorphic on \(\Re t>0\), is positive on
\((0,\infty)\), and extends smoothly from the right to \(t=0\). For every
\(k\ge0\),
\begin{equation}
 S_\mu^{(k)}(t)=(-1)^kk!\int_0^\infty
 \frac{\lambda^k}{(1+t\lambda)^{k+1}}\dd\mu(\lambda).
 \label{eq:derivatives}
\end{equation}
In particular, \(0<S_\mu(t)\le1\) and \(S_\mu'(t)\le0\) for real \(t>0\).
For every integer \(N\ge0\),
\begin{align}
 S_\mu(t)&=\sum_{n=0}^{N-1}(-1)^n m_nt^n+R_N(t),\label{eq:remainder}\\
 R_N(t)&=(-t)^N\int_0^\infty
       \frac{\lambda^N}{1+t\lambda}\dd\mu(\lambda),
 &0\le(-1)^N R_N(t)\le m_Nt^N.\label{eq:remainder-bound}
\end{align}
The empty sum is zero when \(N=0\).
\end{lemma}
\begin{proof}
On the right half-plane, \(|1+t\lambda|\ge1\). Differentiation on compact
sets is justified by the finite moments, giving
\eqref{eq:derivatives}; the same domination gives continuity of every
right derivative at zero. Holomorphy follows from dominated local
holomorphic integration. The signs are immediate. Finally, integrate the
finite geometric identity
\((1+z)^{-1}=\sum_{n<N}(-z)^n+(-z)^N/(1+z)\)
with \(z=t\lambda\ge0\). No infinite geometric expansion has been used.
\end{proof}

\begin{proposition}[Existence and a residual certificate]
\label{prop:inverse}
The map \(F_\mu\) is an increasing bijection of \((0,\infty)\) onto itself.
Equation~\eqref{eq:implicit} has exactly one positive solution for every
\(y>0\). It satisfies
\begin{equation}
 0<T_\mu(y)\le y,
 \qquad T_\mu(y)\ge y(1-m_1y)
 \quad\text{when }m_1y\le1.
 \label{eq:root-bounds}
\end{equation}
For any trial value \(v>0\),
\begin{equation}
 |v-T_\mu(y)|\le|v-yS_\mu(v)|.
 \label{eq:residual}
\end{equation}
The inverse is smooth from the right at zero, with
\begin{equation}
 T_\mu(y)=y-m_1y^2+(m_1^2+m_2)y^3+O(y^4).
 \label{eq:first-inverse}
\end{equation}
\end{proposition}
\begin{proof}
Since
\[
 F_\mu'(t)=\frac{S_\mu(t)-tS_\mu'(t)}{S_\mu(t)^2}\ge1,
\]
\(F_\mu\) is strictly increasing. Moreover, \(F_\mu(t)\to0\) at zero and
\(F_\mu(t)\ge t\to\infty\). The upper root bound follows from \(S\le1\).
The identity \(S(t)\ge1-m_1t\), together with \(T\le y\), gives the lower
bound. For fixed \(y\), the residual function
\(H(t)=t-yS_\mu(t)\) has derivative at least one on the positive axis.
The real mean-value theorem proves \eqref{eq:residual}.

At the boundary point \((t,y)=(0,0)\), the equation has derivative one
in \(t\). Repeated one-sided implicit differentiation gives smoothness
and determines every Taylor coefficient. Equivalently, extend each
finite smooth jet across the endpoint and apply the finite-order
implicit-function theorem. Substitution to degree three yields
\eqref{eq:first-inverse}.
\end{proof}

\begin{remark}[What a residual certificate does not know]
The estimate \eqref{eq:residual} evaluates the actual transform. It does
not turn a formal coefficient list into a selection of \(\mu\).
The information-theoretic lower bound in Section~\ref{sec:transport}
will show why that distinction cannot be omitted.
\end{remark}

\section{Identifiability and the sharp Gevrey threshold}
\label{sec:threshold}

\begin{definition}
A formal series \(\sum a_nz^n\) has \emph{Gevrey order at most \(s\)}
if \(|a_n|\le CA^n(n!)^s\) for some finite \(C,A>0\). It has
\emph{exact order \(s\)} if it has order at most \(s\), but no smaller
order. We use this coefficient convention, not a convention shifted by
one derivative. An assertion about this formal order alone does not
assert a sectorial Borel summation theorem.
\end{definition}

\begin{theorem}[Inverse identifiability equals moment determinacy]
\label{thm:identifiability}
Fix a Stieltjes moment sequence \(m=(m_n)_{n\ge0}\), with \(m_0=1\),
that has at least one representing measure in \(\mathcal P_\infty\).
Every representing measure produces the same formal inverse
\begin{equation}
 \widehat T(y)=y\widehat S(\widehat T(y)),
 \qquad \widehat S(t)=\sum_{n\ge0}(-1)^nm_nt^n.
 \label{eq:formal-implicit}
\end{equation}
Different representing measures produce different actual inverse germs.
Conversely, the formal inverse uniquely determines \(m\). Hence a
unique positive Stieltjes inverse realization exists exactly when the
moment problem is determinate.
\end{theorem}
\begin{proof}
Finite Taylor substitution in Lemma~\ref{lem:regularity} and
Proposition~\ref{prop:inverse} gives \eqref{eq:formal-implicit} to every
finite order. In the formal power-series ring, the equation recursively
determines a unique series beginning with \(y\).

Suppose two actual inverse germs agree. Inverting their increasing
functions shows that their transforms agree on an interval adjacent to
zero. The identity theorem extends equality to the right half-plane.
At \(t=1\), equality of all derivatives implies equality of
\[
 \int\frac{\lambda^k}{(1+\lambda)^{k+1}}\dd\mu(\lambda)
 \quad(k\ge0).
\]
Push \(\mu\) forward by \(u=\lambda/(1+\lambda)\), and multiply the
resulting measure by \(1-u\). We obtain a finite measure on \([0,1]\)
whose polynomial moments are the displayed integrals. Polynomial
approximation of continuous functions on \([0,1]\) gives equality of
these weighted measures. Neither original push-forward has an atom at
\(u=1\), and the weight is strictly positive on \([0,1)\). Division by
the weight, first on \([0,1-\varepsilon]\) and then by monotone
exhaustion, recovers equality of the original measures.

Finally, \(\widehat T(y)=y+O(y^2)\) has a unique formal compositional
inverse \(\widehat F\). The relation
\(\widehat S(t)=t/\widehat F(t)\) reconstructs every moment. This also
shows that no ambiguity is introduced by passing from the transform jet
to the inverse jet.
\end{proof}

\begin{lemma}[A cancellation-free coefficient formula]
\label{lem:coefficients}
Write \(\widehat T(y)=\sum_{j\ge1}b_jy^j\) and
\(c_n=(-1)^n b_{n+1}\). Then \eqref{eq:intro-composition} holds,
\(c_n\ge m_n\), and, for \(s\ge0\), \(C\ge1\),
\begin{equation}
 m_n\le CA^n(n!)^s
 \quad\Longrightarrow\quad
 c_n\le(4AC)^n(n!)^s\quad(n\ge1).
 \label{eq:gevrey-bound}
\end{equation}
Thus the moment sequence and the inverse coefficient sequence have
exactly the same infimum of admissible Gevrey orders; if the order is
attained for either sequence, it is attained for both.
\end{lemma}
\begin{proof}
The formal Lagrange coefficient identity gives
\[
 b_{n+1}=\frac1{n+1}\coef{t}{n}\widehat S(t)^{n+1}.
\]
Every product contributing to degree \(n\) has sign \((-1)^n\), so
there is no cancellation. The \(n+1\) terms with exactly one index equal
to \(n\) contribute \(m_n\) after division by \(n+1\). For the upper
bound, at most \(n\) indices in a weak composition are positive, and
\(m_0=1\). Each product is therefore at most
\(C^nA^n(\prod_jk_j!)^s\le C^nA^n(n!)^s\).
There are \(\binom{2n}{n}\le4^n\) weak compositions of \(n\) into
\(n+1\) parts. Division by \(n+1\) only improves the bound.
\end{proof}

\begin{lemma}[The factorial-square determinacy criterion]
\label{lem:hardy}
A Stieltjes moment sequence satisfying
\(m_n\le CA^n(n!)^2\) is determinate.
\end{lemma}
\begin{proof}
For a representing random variable \(X\ge0\), let \(U\) have the
symmetric law obtained from \(\sqrt X\) by choosing either sign with
probability one half. Its even moments are \(m_n\), and its odd moments
vanish. Choose \(c>0\) with \(c^2A<1\). Monotone convergence gives
\[
 \mathbb E\cosh(cU)
 =\sum_{n\ge0}\frac{c^{2n}m_n}{(2n)!}<\infty,
\]
since \((n!)^2\le(2n)!\). Thus
\(\mathbb E\ee^{c|U|}\le2\mathbb E\cosh(cU)<\infty\).
The moment generating function is holomorphic in the strip
\(|\Re z|<c\), and its Taylor series at zero is fixed by the moments.
Any two representing measures therefore yield identical moment
 generating functions near zero and, by the identity theorem, on the
strip. Their characteristic functions agree on the imaginary axis.
Uniqueness of the Fourier transform of a finite measure gives equality
of the symmetric laws. Squaring recovers the same law of \(X\).
The exponential integrability argument applies to every competitor
because it uses only the shared moments.
\end{proof}

\begin{theorem}[Sharp inverse Gevrey threshold]
\label{thm:gevrey}
In the class of positive Stieltjes inverses, a realizable formal inverse
of Gevrey order at most \(s\le2\) has exactly one realization.
For every real \(s>2\), there are uncountably many distinct realizations
of a single formal inverse of exact Gevrey order \(s\).
\end{theorem}
\begin{proof}
By Lemma~\ref{lem:coefficients}, a Gevrey bound on the inverse implies
the same bound on the moments. For \(s\le2\), this is a factorial-square
bound, so Lemma~\ref{lem:hardy} and
Theorem~\ref{thm:identifiability} give uniqueness. For \(s>2\), take
\(\alpha=1/s\) in Theorem~\ref{thm:gamma} below. That theorem gives a
continuum of distinct measures with common moments of exact order
\(1/\alpha=s\). Lemma~\ref{lem:coefficients} transfers the exact order,
and Theorem~\ref{thm:identifiability} transfers distinctness.
\end{proof}

\begin{remark}[The precise sense of sharpness]
The threshold is a universal \emph{sufficient} growth condition within
a specified positive-measure realization class. It does not say that
every sequence of order greater than two is indeterminate. Nor does it
say that arbitrary smooth or sectorial analytic functions of order at
most two are determined by their positive-real formal expansions.
The classical component is Lemma~\ref{lem:hardy}; compare the treatment
of Hardy's condition and generalized gamma laws in \cite{lin-stoyanov}.
\end{remark}

\section{Generalized gamma measures and an exact exponential defect}
\label{sec:gamma}

Fix parameters
\begin{equation}
 0<\alpha<\tfrac12,\qquad \beta>0,\qquad
 a=\frac\beta\alpha,\quad c=\cos(\pi\alpha),\quad
 A=\sec(\pi\alpha).
 \label{eq:gamma-parameters}
\end{equation}
Define
\begin{align}
 f(\lambda)&=\frac\alpha{\Gamma(a)}\lambda^{\beta-1}
                       \ee^{-\lambda^\alpha},\label{eq:gamma-density}\\
 p_*(\lambda)&=\sin\!\left(\tan(\pi\alpha)\lambda^\alpha-\pi\beta\right),\label{eq:gamma-sine}\\
 \dd\mu_\theta(\lambda)&=f(\lambda)
                \bigl(1-\theta p_*(\lambda)\bigr)\dd\lambda,
 \qquad -1\le\theta\le1.\label{eq:gamma-family}
\end{align}
The notation \(p_*\) is reserved for the bounded sine perturbation;
the unrelated large slope parameter \(p\) will be introduced later.

\begin{theorem}[A moment-null positive family]
\label{thm:gamma}
The measures \(\mu_\theta\) are distinct probabilities with common moments
\begin{equation}
 m_n=\frac{\Gamma((n+\beta)/\alpha)}{\Gamma(\beta/\alpha)}.
 \label{eq:gamma-moments}
\end{equation}
These moments, and their formal inverse coefficients, have exact Gevrey
order \(1/\alpha\).
\end{theorem}
\begin{proof}
The density is nonnegative because \(|p_*|\le1\). Substitution
\(u=\lambda^\alpha\) gives the unperturbed moments. The signed moment
of the perturbation is
\[
 \frac1{\Gamma(a)}\Im\left[
 \ee^{-\ii\pi\beta}
 \int_0^\infty u^{(n+\beta)/\alpha-1}
 \ee^{-(1-\ii\tan(\pi\alpha))u}\dd u\right].
\]
The complex gamma integral is valid because the real part of its
exponential coefficient is positive. Since
\(1-\ii\tan(\pi\alpha)=c^{-1}\ee^{-\ii\pi\alpha}\),
its phase after multiplication by \(\ee^{-\ii\pi\beta}\) is
\(\pi n\). The imaginary part is zero. The case \(n=0\) proves
normalization. The sine perturbation is not zero almost everywhere, so
different parameters give different measures.

Stirling's formula yields
\(\log m_n=\alpha^{-1}n\log n+O(n)\).
The remaining polynomial factors can be absorbed into an exponential
constant, so the order \(1/\alpha\) is attained, while every smaller
order fails. Apply Lemma~\ref{lem:coefficients} to the inverse.
\end{proof}

\begin{theorem}[Closed form for the entire transform defect]
\label{thm:gamma-defect}
For every \(t>0\), the transforms of \eqref{eq:gamma-family} satisfy
\begin{equation}
 S_\theta(t)=S_0(t)+\theta D(t),\qquad
 D(t)=\frac{\pi\alpha}{\Gamma(\beta/\alpha)}
       t^{-\beta}\ee^{-At^{-\alpha}}>0.
 \label{eq:gamma-defect}
\end{equation}
In particular, their common formal expansion omits a nonzero, explicitly
known, flat function.
\end{theorem}
\begin{proof}
Let \(H(t)=\int f(\lambda)p_*(\lambda)/(1+t\lambda)\dd\lambda\).
After \(u=\lambda^\alpha\), introduce
\(v=(1-\ii\tan(\pi\alpha))u\). The phase factors cancel, giving
\begin{equation}
 H(t)=\frac{c^a}{\Gamma(a)}\Im
 \int_{\arg v=-\pi\alpha}
 \frac{v^{a-1}\ee^{-v}}{1-tc^{1/\alpha}v^{1/\alpha}}\dd v.
 \label{eq:gamma-contour}
\end{equation}
Powers are taken with their principal branches in the wedge
\(-\pi\alpha\le\arg v\le0\). The original and transformed integrals
are absolutely convergent.

Rotate the ray to the positive real axis, passing below the simple pole
\(r=t^{-\alpha}/c\). There are no other poles in the open wedge:
the argument of \(v^{1/\alpha}\) lies strictly between \(-\pi\) and
zero. Large arcs vanish exponentially because \(\pi\alpha<\pi/2\),
and small arcs vanish because \(a>0\). The integral along the real axis
is interpreted as its real principal value plus the lower indentation.
The residue is
\begin{equation}
 \Res_{v=r}\frac{v^{a-1}\ee^{-v}}
                        {1-tc^{1/\alpha}v^{1/\alpha}}
 =-\alpha r^a\ee^{-r}.
\end{equation}
The lower indentation, traversed from left to right, contributes
\(\ii\pi\) times the residue. All remaining boundary values are real.
Thus
\[
 H(t)=-\frac{\pi\alpha c^a}{\Gamma(a)}r^a\ee^{-r}
     =-\frac{\pi\alpha}{\Gamma(a)}t^{-\beta}\ee^{-At^{-\alpha}}.
\]
The definition \(1-\theta p_*\) of the density now gives
\eqref{eq:gamma-defect}. Exponential decay dominates every power of
\(t\), proving flatness.
\end{proof}

\begin{example}[A completely explicit Gevrey-three model]
\label{ex:cubic}
For \(\alpha=\beta=1/3\), one has \(m_n=(3n)!\) and
\begin{align}
 S_0(t)&=\int_0^\infty\frac{\ee^{-u}}{1+tu^3}\dd u,
 &D(t)&=\frac\pi3t^{-1/3}\ee^{-2t^{-1/3}}.\label{eq:cubic}
\end{align}
Every \(\theta\in[-1,1]\) gives a positive Stieltjes transform
\(S_0+\theta D\), and the corresponding inverses share the formal series
\begin{equation}
 \widehat T(y)=y-6y^2+756y^3-376056y^4
                  +488904336y^5+\cdots.
 \label{eq:cubic-inverse}
\end{equation}
The continuum of actual inverse functions cannot be distinguished by
computing additional coefficients of this series.
\end{example}

\begin{remark}[The endpoint and the classical dichotomy]
For the unperturbed generalized gamma law with arbitrary \(\alpha>0\),
\eqref{eq:gamma-moments} has order \(1/\alpha\). Thus
\(\alpha\ge1/2\) falls within the determinacy criterion, whereas the
family above proves indeterminacy for \(\alpha<1/2\). This recovers a
classical generalized gamma dichotomy rather than replacing it with a
new moment theorem. The contour formula makes an additional quantitative
feature visible: its action \(A=\sec(\pi\alpha)\) diverges as the
boundary is approached from below.
\end{remark}

\section{Bilateral geometric measures and theta-scale ambiguity}
\label{sec:qfamily}

Fix \(0<q<1\), put \(h=-\log q\), and set
\begin{equation}
 w_n=q^{n(n-1)/2},\qquad \lambda_n=q^{-n}
 \quad(n\in\Z).
\end{equation}
The Gaussian decay in \(n\) makes every sum used below absolutely
convergent. Define
\begin{equation}
 M=\sum_{n\text{ even}}w_n=\sum_{n\text{ odd}}w_n,
 \qquad
 \mu_+=\frac1M\sum_{n\text{ even}}w_n\delta_{q^{-n}},\quad
 \mu_-=\frac1M\sum_{n\text{ odd}}w_n\delta_{q^{-n}}.
 \label{eq:q-measures}
\end{equation}
The equality defining \(M\) follows by pairing \(n\) with \(1-n\).
Write \(S_\pm=S_{\mu_\pm}\), and
\(S_\eta=(1-\eta)S_-+\eta S_+\) for \(0\le\eta\le1\).
These discrete measures belong to the classical Stieltjes--Wigert
moment setting \cite{chihara,christiansen,christiansen-koelink}.

\begin{proposition}[Common moments and exact difference equations]
\label{prop:q-moments}
Both measures in \eqref{eq:q-measures} have moments
\begin{equation}
 m_k=q^{-k(k+1)/2}\quad(k\ge0).
 \label{eq:q-moments}
\end{equation}
Their transforms obey
\begin{equation}
 S_+(qt)=1-tS_-(t),\qquad S_-(qt)=1-tS_+(t).
 \label{eq:coupled-q}
\end{equation}
Consequently every \(S_\eta\) satisfies the same exact equation
\begin{equation}
 S_\eta(q^2t)=1-qt+qt^2S_\eta(t).
 \label{eq:q2}
\end{equation}
The common formal series satisfies
\(\widehat S(qt)=1-t\widehat S(t)\), although not every member of this
particular affine family satisfies that first-order equation exactly.
\end{proposition}
\begin{proof}
The index shift \(n=m+k\) gives
\(w_nq^{-kn}=q^{-k(k+1)/2}w_m\).
It preserves or reverses parity according to \(k\); either parity has
normalizing sum \(M\). This proves the moments. For the difference
equations, shift \(n=m+1\), use \(w_{m+1}=q^mw_m\), and use
\[
 \frac{q^m}{1+tq^{-m}}=q^m-\frac{t}{1+tq^{-m}}.
\]
The shifted constant sum is again \(M\), giving
\eqref{eq:coupled-q}. A second iteration proves \eqref{eq:q2}.
The formal first-order identity follows directly from
\eqref{eq:q-moments} by comparing coefficients.
\end{proof}

We use the following normalization of the theta function:
\begin{align}
 \Theta_q(z)&=\sum_{n\in\Z}q^{n(n-1)/2}z^n,\label{eq:theta}\\
 (a;q)_\infty&=\prod_{j=0}^\infty(1-aq^j),\nonumber\\
 \Theta_q(z)&=(-z;q)_\infty(-q/z;q)_\infty(q;q)_\infty.
 \label{eq:triple}
\end{align}
Equation~\eqref{eq:triple} is Jacobi's triple-product identity
\cite{dlmf}. It also fixes all normalization constants in what follows.

\begin{theorem}[Exact reciprocal-theta defect]
\label{thm:theta-defect}
For every \(t>0\),
\begin{equation}
 D(t):=S_+(t)-S_-(t)
 =\frac{(q;q)_\infty^3}{M\Theta_q(t)}
 =\frac{(q;q)_\infty^2}
 {M(-t;q)_\infty(-q/t;q)_\infty}>0.
 \label{eq:theta-defect}
\end{equation}
In particular,
\begin{equation}
 D(qt)=tD(t).
 \label{eq:defect-q}
\end{equation}
\end{theorem}
\begin{proof}
Let
\[
 G(t)=\sum_{n\in\Z}\frac{(-1)^nw_n}{1+tq^{-n}}=MD(t).
\]
The series is normally convergent on compact subsets of
\(\C^*\setminus\{-q^n:n\in\Z\}\). The same index shift as above gives
\(G(qt)=tG(t)\), because \(\sum(-1)^nq^nw_n=0\).
Also \(\Theta_q(qt)=t^{-1}\Theta_q(t)\).

The product \(G(t)\Theta_q(t)\) is holomorphic on \(\C^*\): the zeros
in the triple product cancel the simple poles of \(G\). It is invariant
under multiplication by \(q\). It is bounded on a closed fundamental
annulus and hence, by this invariance, on all of \(\C^*\). Its
singularity at zero is removable. Liouville's theorem makes it constant.
At \(t=-1\), the residue of \(G\) is one, whereas
\(\Theta_q(t)/(1+t)\to(q;q)_\infty^3\).
The constant is therefore \((q;q)_\infty^3\), proving the formula.
Every factor in the product form is positive for \(t>0\).
\end{proof}

\subsection{A precise log-periodic amplitude}
For real \(u\), define
\begin{equation}
 P_h(u)=\sum_{n\in\Z}\ee^{-h(n+u)^2/2},\qquad
 K_q=\frac{(q;q)_\infty^3}{M},\qquad
 x=\frac{\log(1/t)}h.
 \label{eq:periodic}
\end{equation}
Completing the square in the theta series yields the \emph{exact}
identity
\begin{equation}
 D(t)=K_q\,
 \frac{\exp[-\frac h2(x-\frac12)^2]}{P_h(x-\frac12)}.
 \label{eq:theta-gaussian}
\end{equation}
The function \(P_h\) is positive, smooth, and one-periodic. Its Fourier
coefficients follow by integrating the periodized Gaussian over one
period:
\begin{equation}
 P_h(u)=\sqrt{\frac{2\pi}{h}}
 \sum_{k\in\Z}\ee^{-2\pi^2k^2/h}\ee^{2\pi\ii ku}.
 \label{eq:poisson}
\end{equation}
The Gaussian Fourier transform justifies this absolutely convergent
Fourier expansion; in particular, \(P_h\) is not constant.

It follows that \(D\) is flat, but
\begin{equation}
 \frac{D(t)}{\exp[-\frac h2(x-\frac12)^2]}
 \label{eq:nonlimit}
\end{equation}
has no limit as \(t\downarrow0\). Indeed, choose two phases at which
\(P_h\) differs and let \(x\) tend to infinity along those phases.
Thus a single real constant-amplitude monomial at this scale is
insufficient. A periodic coefficient, or an equivalent oscillatory
extension, is needed. This observation does not exclude every possible
extended transseries representation.

\begin{theorem}[The ambiguity matches the least moment bound]
\label{thm:least}
Let
\begin{equation}
 E_*(t)=\min_{N\in\N} t^Nq^{-N(N+1)/2},
 \qquad \N=\{0,1,2,\ldots\}.
 \label{eq:least-bound}
\end{equation}
For \(x\ge1/2\),
\begin{align}
 E_*(t)&=\exp\!\left[-\frac h2(x-\tfrac12)^2
              +\frac h2\dist(x-\tfrac12,\Z)^2\right],
 \label{eq:least-exact}\\
 \frac{D(t)}{E_*(t)}
 &=\frac{K_q\exp[-\frac h2\dist(x-\tfrac12,\Z)^2]}
           {P_h(x-\tfrac12)}.\label{eq:least-ratio}
\end{align}
In particular, \(D(t)\asymp_q E_*(t)\) as \(t\downarrow0\): the ratio
is bounded above and below by strictly positive constants depending
only on \(q\).
\end{theorem}
\begin{proof}
The logarithm of the \(N\)-th bound is
\[
 \log(t^Nm_N)=\frac h2\bigl(N-(x-\tfrac12)\bigr)^2
                      -\frac h2(x-\tfrac12)^2.
\]
For \(x\ge1/2\), a nearest integer to \(x-1/2\) is nonnegative. This
proves \eqref{eq:least-exact}. Divide \eqref{eq:theta-gaussian} by it.
The periodic continuous denominator has a positive minimum and finite
maximum, and \(0\le\dist(u,\Z)\le1/2\).
\end{proof}

\begin{remark}[A bound, not an equality of true remainders]
The quantity \(E_*\) is the best member of the explicit family of
bounds in \eqref{eq:remainder-bound}. The theorem does not assert that
it equals the smallest actual truncation error of either transform.
It does show that no coefficient-only reconstruction can generally
remove ambiguity at that scale. Also,
\(\log m_n=\frac h2n^2+O(n)\), so this example lies beyond every finite
Gevrey order.
\end{remark}

\subsection{Even the exact first-order equation can remain nonunique}
\label{subsec:shifted}
One might try to remove the ambiguity by imposing the first-order
\(q\)-Euler equation exactly. Within the affine family \(S_\eta\), it
selects \(\eta=1/2\). In the full positive-measure class, that conclusion
is false.

\begin{proposition}[Shifted lattices with the same exact equation]
\label{prop:shifted}
For \(r>0\), define
\begin{equation}
 Z_r=\sum_{n\in\Z}r^{-n}w_n=\Theta_q(1/r),\qquad
 \nu_r=\frac1{Z_r}\sum_{n\in\Z}r^{-n}w_n\delta_{r q^{-n}}.
 \label{eq:shifted-lattice}
\end{equation}
All these probabilities have moments \eqref{eq:q-moments}. Their
transforms \(U_r=S_{\nu_r}\) all satisfy
\begin{equation}
 U_r(qt)=1-tU_r(t).
 \label{eq:exact-q-euler}
\end{equation}
If \(r/s\notin q^\Z\), the transforms and their inverse germs are
distinct. Their difference changes sign infinitely often at zero.
The corresponding inverse germs also cross infinitely often.
\end{proposition}
\begin{proof}
For the moments, shift \(n=m+k\) as in
Proposition~\ref{prop:q-moments}; the factors \(r^k\) from the support
and \(r^{-k}\) from the weight cancel. If
\(a_n=Z_r^{-1}r^{-n}w_n\), then
\(a_{n+1}=(rq^{-n})^{-1}a_n\). Shifting the transform sum gives
\[
 U_r(qt)=\sum_n\frac{a_n(rq^{-n})^{-1}}{1+trq^{-n}}
        =\sum_na_n(rq^{-n})^{-1}-tU_r(t)=1-tU_r(t).
\]
The first sum is \(\sum a_{n+1}=1\).

When \(r/s\notin q^\Z\), the supports are disjoint, so the measures
are different. By Theorem~\ref{thm:identifiability}, so are their
transforms. For \(V=U_r-U_s\), equation~\eqref{eq:exact-q-euler} gives
\(V(qt)=-tV(t)\). Choose \(t_*>0\) with \(V(t_*)\ne0\). Iteration
shows alternating signs at \(q^nt_*\); continuity gives a zero between
successive such points.

For inverse transport, write \(y=t/U_s(t)\), so \(T_s(y)=t\).
The residual for the other inverse is
\(t-yU_r(t)=-yV(t)\). Since its residual is strictly increasing,
\(\operatorname{sgn}(T_r(y)-T_s(y))=\operatorname{sgn}V(t)\).
The sequence \(t=q^nt_*\) tends to zero and so does \(y\); its signs
alternate. This proves infinitely many inverse crossings.
\end{proof}

\begin{corollary}[A simultaneous germ-field obstruction]
The two inverse germs in Proposition~\ref{prop:shifted} cannot both lie
in a field of actual real-valued germs at zero. In particular, after
\(y=1/x\), they cannot both belong to one Hardy field at infinity.
\end{corollary}
\begin{proof}
Their difference is a nonzero germ but has zeros arbitrarily close to
zero. It therefore has no reciprocal germ. A field cannot contain such
a nonzero element.
\end{proof}

\begin{remark}
This is a statement about simultaneous realization, not a proof that
either inverse individually lies in no Hardy field. The basic
moment-indeterminate lattice families are classical. Their role here is
to show exactly which extra analytic equations still fail to turn formal
data into unique inverse data.
\end{remark}

\section{Nonlinear transport and an unavoidable error floor}
\label{sec:transport}

In this section \(S_a=S_{a_0}+(a-a_0)D\) denotes either the gamma family
with physical interval \([-1,1]\), or the ordered theta family with
physical interval \([0,1]\). The word \emph{physical} means simply that
the associated measure is positive; it does not refer to an external
physical model. In both families \(D(t)>0\), and every member has the
same moments.

\begin{lemma}[Exact derivative with respect to the hidden parameter]
\label{lem:parameter-derivative}
For \(y>0\) and a physical parameter \(a\),
\begin{equation}
 \partial_aT_a(y)=
 \frac{yD(T_a(y))}{1-yS_a'(T_a(y))}>0.
 \label{eq:parameter-derivative}
\end{equation}
For two physical parameters \(a,b\),
\begin{equation}
 T_a(y)-T_b(y)
 =\int_b^a\frac{yD(T_\theta(y))}
                   {1-yS_\theta'(T_\theta(y))}\dd\theta.
 \label{eq:gap-integral}
\end{equation}
Every finite Taylor expansion of \(T_a\) at zero is independent of \(a\).
\end{lemma}
\begin{proof}
Differentiate \(T_a=yS_a(T_a)\). The denominator is at least one by
positivity of the measure, proving the formula and its sign. Integration
gives the second statement. Parameter endpoints are handled by one-sided
derivatives or by local analytic continuation. Shared Taylor
coefficients follow from Theorem~\ref{thm:identifiability}.
\end{proof}

\begin{theorem}[Two quantitative inverse-gap laws]
\label{thm:gap-laws}
The following expansions hold uniformly for distinct physical
parameters \(a,b\), including in the limit \(a-b\to0\).

For the gamma family, with the parameters of \eqref{eq:gamma-parameters},
\begin{equation}
 \frac{T_a(y)-T_b(y)}{(a-b)yD(y)}
 =1-\alpha A m_1y^{1-\alpha}
       +(\beta-1)m_1y+O(y^{2-2\alpha}).
 \label{eq:gamma-gap}
\end{equation}
For the theta family, define the bounded one-periodic function
\begin{equation}
 a_q(x)=-\frac12+
 \frac{P_h'(x-\frac12)}{hP_h(x-\frac12)}.
 \label{eq:aq}
\end{equation}
With \(x=\log(1/y)/h\) and \(m_1=q^{-1}\),
\begin{equation}
 \frac{T_a(y)-T_b(y)}{(a-b)yD(y)}
 =1-m_1y\bigl(x+a_q(x)+1\bigr)
       +O\!\left(y^2(1+x)^2\right).
 \label{eq:q-gap}
\end{equation}
\end{theorem}
\begin{proof}
The shared finite moments give, uniformly in the physical parameter,
\[
 T_a(y)=y-m_1y^2+O(y^3),\qquad
 (1-yS_a'(T_a(y)))^{-1}=1-m_1y+O(y^2).
\]
Uniformity follows either from bounded densities relative to a common
base in the gamma case, or from convexity in the theta case, and the
finite-moment derivative bounds.

For gamma, expand the logarithm of the exact defect:
\begin{align*}
 \log\frac{D(T_a(y))}{D(y)}
 &=-\beta\log\frac{T_a(y)}y
   -Ay^{-\alpha}\left[\left(\frac{T_a(y)}y\right)^{-\alpha}-1\right]\\
 &=-\alpha A m_1y^{1-\alpha}+\beta m_1y+O(y^{2-\alpha}).
\end{align*}
Exponentiating introduces an error of order \(y^{2-2\alpha}\).
Multiplying by the denominator expansion gives the integrand in
\eqref{eq:gap-integral}, divided by \(yD(y)\). The asserted gamma law
follows. Because \(\alpha<1/2\), the error has order strictly greater
than one, so the displayed linear term is informative.

For theta, direct differentiation of \eqref{eq:theta-gaussian} gives
\begin{equation}
 t\frac{D'(t)}{D(t)}=x+a_q(x),
 \qquad x=\frac{\log(1/t)}h.
 \label{eq:theta-slope}
\end{equation}
Since \(a_q\) and all its real derivatives are bounded, integration over
\(\log(T_a(y)/y)=-m_1y+O(y^2)\) yields
\[
 \log\frac{D(T_a(y))}{D(y)}
 =-m_1y(x+a_q(x))+O(y^2(1+x)).
\]
Exponentiate, multiply by the denominator factor, and use
\eqref{eq:gap-integral}. All estimates are uniform in \(a\), hence also
in the averaged interval from \(b\) to \(a\).
\end{proof}

For the cubic example, equation~\eqref{eq:gamma-gap} becomes particularly
simple:
\begin{equation}
 \frac{T_a(y)-T_b(y)}{(a-b)yD(y)}
 =1-4y^{2/3}-4y+O(y^{4/3}).
 \label{eq:cubic-gap}
\end{equation}
Thus a flat perturbation does not merely persist: its first nonlinear
relative correction can be computed without resolving the full
coefficient expansion.

\begin{theorem}[A coefficient-only information floor]
\label{thm:floor}
Let an estimator \(\mathcal A(m,y)\) depend only on \(y\) and the full
common moment sequence, or equivalently on the full common inverse
coefficient sequence. For any two physical members,
\begin{equation}
 \max\{|\mathcal A-T_a(y)|,|\mathcal A-T_b(y)|\}
 \ge\frac12|T_a(y)-T_b(y)|.
 \label{eq:information-floor}
\end{equation}
For the theta endpoints, this lower bound is asymptotic to
\(yD(y)/2\) and hence is bounded below by \(c_qyE_*(y)\) for all
sufficiently small \(y\), with \(c_q>0\). For the gamma endpoints
\(a=1,b=-1\), it is asymptotic to \(yD(y)\).
\end{theorem}
\begin{proof}
The estimator receives identical input for both members. The triangle
inequality gives \eqref{eq:information-floor}. Apply
Theorem~\ref{thm:gap-laws}, and for theta use Theorem~\ref{thm:least}.
\end{proof}

\begin{remark}
This lower bound concerns a worst-case error over a family of admissible
realizations. It is not a limit on methods supplied with the actual
measure, an exact integral, a summation prescription, or another datum
that fixes the hidden parameter. In the gamma case, we do not identify
the defect action \(A\) with the least-moment-bound action: they are
different quantities.
\end{remark}

\section{Convergent expansions in the flat ambiguity parameter}
\label{sec:sectors}

A finite truncation of the common formal power series cannot describe
the ambiguity in Section~\ref{sec:transport}. A different expansion can:
keep one actual realization as the baseline, and expand in the hidden
parameter. In this section the letters \(a,a_0\) denote family
parameters, not the gamma quotient \(\beta/\alpha\).

Fix a physical baseline \(a_0\), a small \(y>0\), and set
\begin{equation}
 t_0=T_{a_0}(y),\quad S=S_{a_0},\quad
 H(t)=t-yS(t),\quad B_0=1-yS'(t_0),\quad D_0=D(t_0).
 \label{eq:sector-data}
\end{equation}
The deformed inverse equation is
\begin{equation}
 H(t)=\delta yD(t),\qquad \delta=a-a_0.
 \label{eq:deformed-equation}
\end{equation}
Complex \(\delta\) is allowed in the local analytic discussion, even
though only a bounded real interval corresponds to positive measures.

\begin{theorem}[Exact sector coefficients and an explicit tail bound]
\label{thm:sectors}
Assume \(ym_1\le1/2\), and choose \(0<r<t_0/2\). On a neighborhood
of the closed disc \(|t-t_0|\le r\), define the removable analytic quotient
\begin{equation}
 A_y(t)=\frac{t-t_0}{H(t)},\quad A_y(t_0)=B_0^{-1},\qquad
 \phi_y(t)=yD(t)A_y(t).
 \label{eq:phi}
\end{equation}
Let \(M_D\) be any bound for \(|D|\) on the disc and set \(M=2yM_D\).
For \(\rho=|\delta|M/r<1\), the branch through \(t_0\) is represented by
an absolutely convergent series
\begin{align}
 T_{a_0+\delta}(y)&=t_0+\sum_{k\ge1}\delta^k A_k(y),\label{eq:sectors}\\
 A_k(y)&=\frac1{k!}
  \left.\frac{\mathrm d^{k-1}}{\mathrm dt^{k-1}}
  \bigl(\phi_y(t)^k\bigr)\right|_{t=t_0}.\label{eq:sector-coefficients}
\end{align}
Its coefficients and tails satisfy
\begin{align}
 |\delta^kA_k(y)|&\le\frac r k\rho^k,\label{eq:sector-bound}\\
 \left|T_{a_0+\delta}(y)-t_0-\sum_{k=1}^K\delta^kA_k(y)\right|
 &\le\frac{r\rho^{K+1}}{(K+1)(1-\rho)}.\label{eq:sector-tail}
\end{align}
For physical \(a_0+\delta\), this branch is the positive inverse already
defined in Section~\ref{sec:framework}.
\end{theorem}
\begin{proof}
The disc lies in the right half-plane. There \(|S'(t)|\le m_1\), by
\eqref{eq:derivatives}. Since \(H(t_0)=0\),
\[
 \frac{H(t)}{t-t_0}
 =1-y\int_0^1 S'(t_0+s(t-t_0))\dd s.
\]
The right side differs from one by at most \(ym_1\le1/2\). Thus it is
nonzero and \(|A_y|\le2\) throughout the disc. Equation
\eqref{eq:deformed-equation} is equivalent there to
\(t-t_0=\delta\phi_y(t)\).

On the boundary, \(|\delta\phi_y|\le|\delta|M<r=|t-t_0|\).
Rouch\'e's theorem gives one root, counted with multiplicity, in the
disc. It is simple and depends holomorphically on \(\delta\) for
\(|\delta|<r/M\). For real physical parameters this root is real by
conjugation symmetry, and positive because \(r<t_0\). Uniqueness of
the positive inverse identifies it with \(T_{a_0+\delta}\).

The analytic Lagrange formula gives \eqref{eq:sector-coefficients}.
For completeness, putting \(w=t-t_0\), the coefficient is
\(k^{-1}[w^{k-1}]\phi_y(t_0+w)^k\).
Cauchy's coefficient estimate on \(|w|=r\) gives
\(|A_k|\le M^k/(kr^{k-1})\). Summing these bounds and using
\(1/k\le1/(K+1)\) for \(k\ge K+1\) proves
\eqref{eq:sector-tail}.
\end{proof}

\begin{proposition}[The first two exact sectors]
\label{prop:first-sectors}
With the data of \eqref{eq:sector-data},
\begin{align}
 A_1(y)&=\frac{yD_0}{B_0},\label{eq:a1}\\
 A_2(y)&=\frac{y^2D_0D'(t_0)}{B_0^2}
       +\frac{y^3D_0^2S''(t_0)}{2B_0^3}.\label{eq:a2}
\end{align}
\end{proposition}
\begin{proof}
Expand \(H(t_0+w)=B_0w-yS''(t_0)w^2/2+O(w^3)\). Hence
\(A_y(t_0)=B_0^{-1}\) and
\(A_y'(t_0)=yS''(t_0)/(2B_0^2)\).
Substitution in \eqref{eq:sector-coefficients}, or direct implicit
differentiation twice, gives the formulas.
\end{proof}

The coefficients contain successively higher powers of \(D_0\),
multiplied by slope factors. Thus \eqref{eq:sectors} resolves the
beyond-all-orders ambiguity rather than merely recomputing the common
power series. It is an expansion \emph{around an actual baseline}; it
does not supply a canonical choice of that baseline from its formal jet.

\subsection{Explicit discs in the two families}
The next bounds give completely specified sufficient convergence
conditions. They are deliberately simple, rather than asymptotically
optimal. Section~\ref{sec:fold} will determine the leading optimal
radius.

\begin{proposition}[A gamma disc]
\label{prop:gamma-disc}
For the gamma defect, put
\begin{equation}
 p=\alpha A t_0^{-\alpha},\qquad
 r=\frac{t_0}{8(1+\beta+p)}.
 \label{eq:gamma-disc}
\end{equation}
Then \(M_D=\ee^{3/8}D_0\) is valid in
Theorem~\ref{thm:sectors}. For all sufficiently small \(y\),
\begin{equation}
 \frac Mr\le32\ee^{3/8}(1+\beta+p)D_0\longrightarrow0.
 \label{eq:gamma-disc-bound}
\end{equation}
\end{proposition}
\begin{proof}
On \(|t-t_0|\le t_0/2\), using the principal powers,
\[
 \left|\frac{D'(t)}{D(t)}\right|
 \le\frac{2\beta+2^{1+\alpha}p}{t_0}
 \le\frac{2\beta+3p}{t_0}.
\]
The nonvanishing explicit defect has a holomorphic logarithm on this
disc. Integrating its logarithmic derivative along a straight segment
of length at most \(r\) bounds the logarithmic variation by \(3/8\).
This proves the bound for \(M_D\). Since \(t_0/y\to1\), eventually
\(y/t_0\le2\); substituting \(M=2yM_D\) gives
\eqref{eq:gamma-disc-bound}. Its limit is zero by exponential flatness.
\end{proof}

\begin{proposition}[A theta disc]
\label{prop:q-disc}
For the theta defect, let \(t_0<1\), and put
\begin{equation}
 x_0=\frac{\log(1/t_0)}h,\qquad
 B=x_0+\frac4{1-q},\qquad r=\frac{t_0}{8B}.
 \label{eq:q-disc}
\end{equation}
Then \(M_D=\ee^{1/4}D_0\) is valid in
Theorem~\ref{thm:sectors}, and eventually
\begin{equation}
 \frac Mr\le32\ee^{1/4}BD_0\longrightarrow0.
 \label{eq:q-disc-bound}
\end{equation}
\end{proposition}
\begin{proof}
Logarithmically differentiate the product in
\eqref{eq:theta-defect}:
\[
 t\frac{D'(t)}{D(t)}
 =-\sum_{n\ge0}\frac{tq^n}{1+tq^n}
   +\sum_{n\ge0}\frac{q^{n+1}}{t+q^{n+1}}.
\]
For \(|t-t_0|\le t_0/2\), the first absolute sum is at most
\(3/[2(1-q)]\). In the second, the terms with \(n+1\le x_0\) are each
at most one, while the remaining sum is at most \(2/(1-q)\).
Thus \(|tD'/D|\le B\), and \(|D'/D|\le2B/t_0\).
No zero or pole of \(D\) lies in this disc. Integrating its logarithmic
derivative over radius \(r\) gives variation at most \(1/4\). The last
estimate follows from \(y/t_0\le2\) and the log-Gaussian flatness of
\(D_0\).
\end{proof}

Both propositions imply convergence uniformly over every bounded
physical parameter interval when \(y\) is sufficiently small. Their
bounds are also effective once the actual baseline and its defect are
numerically enclosed. The floating-point examples later in the article
do not by themselves provide such outward-rounded enclosures.

\section{A universal Lambert--\texorpdfstring{\(W\)}{W} limit and the parameter radius}
\label{sec:fold}

The scale of the sharp parameter radius is larger than the elementary
Cauchy-disc radius by a bounded factor. More importantly, it has a
universal limit. Continue to fix a physical baseline \(a_0\), and use
\eqref{eq:sector-data}. Define
\begin{equation}
 p=\begin{cases}
    \alpha A t_0^{-\alpha},&\text{gamma},\\
    \log(1/t_0)/h,&\text{theta},
   \end{cases}
 \qquad
 v=\frac{p(t-t_0)}{t_0},\qquad
 \kappa=\delta\frac{pyD_0}{t_0}.
 \label{eq:fold-scaling}
\end{equation}
In both cases \(p\to\infty\) as \(y\downarrow0\). Equation
\eqref{eq:deformed-equation} becomes
\begin{equation}
 \kappa=\Psi_y(v):=
 \frac{pH(t_0(1+v/p))}{t_0}
 \frac{D_0}{D(t_0(1+v/p))}.
 \label{eq:psi}
\end{equation}
On every fixed bounded complex \(v\)-set, this formula is holomorphic
for all sufficiently small \(y\). It no longer involves a subtraction
of two unspecified formal series.

\begin{lemma}[Uniform local convergence of the deformation map]
\label{lem:psi-limit}
For either family,
\begin{equation}
 \Psi_y(v)\longrightarrow v\ee^{-v}
 \label{eq:psi-limit}
\end{equation}
locally uniformly in \(v\in\C\), with convergence of every derivative.
\end{lemma}
\begin{proof}
Since \(H(t_0)=0\) and \(|S'|\le m_1\) in the right half-plane,
\[
 \frac{pH(t_0(1+v/p))}{t_0}=v+O_R(y)
 \quad(|v|\le R).
\]
For gamma, the explicit logarithm gives
\[
 \log\frac{D(t_0(1+v/p))}{D_0}
 =-\beta\log(1+v/p)
  -\frac p\alpha\bigl((1+v/p)^{-\alpha}-1\bigr)
 =v+O_R(p^{-1}).
\]
For theta, integrate \eqref{eq:theta-slope} over
\(\log(1+v/p)\). The periodic factor and its derivatives are uniformly
bounded in a fixed complex strip about the real phase axis: positivity
on that axis and compactness of one period supply a zero-free strip.
The resulting logarithm is again \(v+O_R(p^{-1})\).
This proves \eqref{eq:psi-limit}. Local uniform holomorphic convergence
implies convergence of all derivatives by Cauchy's formula.
\end{proof}

\begin{theorem}[Cayley coefficients, a real fold, and a sharp leading radius]
\label{thm:radius}
Let \(R_y\) be the radius of the Taylor series in \(\delta\) for the
branch \(T_{a_0+\delta}(y)\) through \(t_0\). Then
\begin{equation}
 R_y\sim\frac{t_0}{\ee\,p\,yD_0}.
 \label{eq:radius}
\end{equation}
In scaled coordinates, the branch converges locally uniformly for
\(|\kappa|<1/\ee\) to
\begin{equation}
 v_0(\kappa)=-W_0(-\kappa)
       =\sum_{k\ge1}\frac{k^{k-1}}{k!}\kappa^k,
 \qquad v_0\ee^{-v_0}=\kappa.
 \label{eq:lambert-limit}
\end{equation}
Consequently, for every fixed \(k\ge1\),
\begin{equation}
 A_k(y)=\frac{t_0}{p}
   \left(\frac{pyD_0}{t_0}\right)^k
   \left(\frac{k^{k-1}}{k!}+O_k(p^{-1})\right).
 \label{eq:cayley-limit}
\end{equation}
There is a real critical point \(v_c(y)\to1\) of \(\Psi_y\), with
critical value \(\kappa_c(y)\to1/\ee\), on the continuation of the
branch from the origin. It is a nondegenerate square-root fold.
\end{theorem}
\begin{proof}
We separate the lower radius bound from the upper bound; the latter
cannot be inferred merely from coefficient limits at fixed \(k\).

\emph{Analytic continuation inside every smaller disc.}
Fix \(r<1/\ee\) and choose \(r<r_1<1/\ee\). The inverse
\(v_0=-W_0(-\kappa)\), defined equivalently by its Taylor series, is
holomorphic on a neighborhood of \(|\kappa|\le r_1\). Its image is
compact, and \((v\ee^{-v})'=\ee^{-v}(1-v)\) is nonzero there.
Uniform local inverse neighborhoods can therefore be chosen around
\(v_0(\kappa)\), with one common sufficiently small radius on the
compact \(\kappa\)-disc. Taylor's theorem and the positive minimum of
the derivative give a uniform nonzero lower bound for
\(|v\ee^{-v}-\kappa|\) on the boundaries of those neighborhoods.
Lemma~\ref{lem:psi-limit} and Rouch\'e's theorem then give a unique nearby
root of \(\Psi_y(v)=\kappa\), analytic in \(\kappa\). Local uniqueness
makes these roots agree on overlapping neighborhoods, producing a
single analytic branch on \(|\kappa|\le r\).
The roots converge uniformly to \(v_0\). The same argument, using the
\(O(p^{-1})\) error in the proof of the lemma, gives that convergence
rate on every smaller fixed disc. Hence the scaled Taylor radius has
liminf at least \(1/\ee\).

\emph{A singularity on the same branch.}
The derivative of \(v\ee^{-v}\) has a simple zero at \(v=1\), and its
second derivative there is \(-1/\ee\). Holomorphic convergence gives
one simple zero \(v_c(y)\) of \(\Psi_y'\) near one. Reality of the
functions implies that this zero is real. On every interval
\([0,1-\varepsilon]\), \(\Psi_y'\) is positive for sufficiently small
\(y\). In a fixed small neighborhood of one, \(\Psi_y''\) is negative,
so \(\Psi_y'\) decreases through its unique zero. It follows that the
branch from zero increases to \(v_c\), at
\(\kappa_c=\Psi_y(v_c)\to1/\ee\).

The Taylor expansion at the critical point is
\(\Psi_y(v)=\kappa_c+\Psi_y''(v_c)(v-v_c)^2/2+O((v-v_c)^3)\).
Since \(\Psi_y''(v_c)\ne0\), the inverse has a genuine square-root
singularity there. An analytic continuation of the Taylor series across
that point on this branch is impossible. Thus the scaled radius has
limsup at most \(1/\ee\). Combining both bounds and undoing
\eqref{eq:fold-scaling} proves \eqref{eq:radius}.

Finally, the inverse of \(v\mapsto v\ee^{-v}\) has the coefficients in
\eqref{eq:lambert-limit} by Lagrange inversion. Uniform convergence of
the scaled analytic inverses on a fixed circle gives convergence of
coefficients, with error \(O_k(p^{-1})\). Rescaling yields
\eqref{eq:cayley-limit}.
\end{proof}

\begin{remark}[What the radius theorem does and does not locate]
The theorem determines the leading asymptotic Taylor radius. It does
not assert that the real fold is the exact nearest complex singularity
for every fixed nonzero \(y\). Competing singularities could alter
subleading information about that radius. The fold location itself,
which is a local real statement, is calculated more precisely next.
For small \(y\), its \(\delta\)-value tends to infinity, so it lies far
outside any fixed positive-measure parameter interval. No loss of
existence of the physical positive inverse is being asserted.
\end{remark}

\subsection{The first fold correction and the surviving theta phase}

\begin{lemma}[A two-term scaled defect expansion]
\label{lem:Q}
Uniformly on bounded complex \(v\)-sets,
\begin{equation}
 \log\frac{D(t_0(1+v/p))}{D_0}
 =v+\frac{Q_y(v)}p+O_R(p^{-2}),
 \label{eq:Qexp}
\end{equation}
where
\begin{equation}
 Q_y(v)=
 \begin{cases}
 -\beta v-\dfrac{1+\alpha}{2}v^2,&\text{gamma},\\[3pt]
 a_q(p)v-\dfrac12v^2,&\text{theta}.
 \end{cases}
 \label{eq:Q}
\end{equation}
Moreover,
\begin{equation}
 \Psi_y(v)=v\ee^{-v}\left(1-\frac{Q_y(v)}p\right)
                      +O_R(p^{-2}),
 \label{eq:psicorrection}
\end{equation}
with the same control for finitely many derivatives on smaller compact
sets.
\end{lemma}
\begin{proof}
The gamma formula follows by taking two further terms in the binomial
and logarithmic expansions used in Lemma~\ref{lem:psi-limit}.
For theta, put \(s=\log(t/t_0)\), so the phase is \(p-s/h\). Integrate
\[
 \frac{\mathrm d}{\mathrm ds}\log D(t_0\ee^s)
 =p-s/h+a_q(p-s/h)
\]
from zero to \(\log(1+v/p)\). The terms of orders one and \(p^{-1}\)
are exactly those in \eqref{eq:Q}. Bounded complex-strip derivatives of
\(a_q\) control the remainder uniformly in its phase.

The residual prefactor in \eqref{eq:psi} is \(v+O_R(y)\). In the gamma
case \(y\asymp t_0\asymp p^{-1/\alpha}=o(p^{-2})\), because
\(\alpha<1/2\). In the theta case \(y\asymp\ee^{-hp}=o(p^{-2})\).
The residual correction is therefore smaller than the stated error.
Exponentiate \eqref{eq:Qexp} with a minus sign and multiply by this
prefactor. Holomorphy again gives derivative control.
\end{proof}

\begin{theorem}[Explicit gamma and theta fold corrections]
\label{thm:fold-correction}
The real critical point and value from Theorem~\ref{thm:radius} satisfy
\begin{equation}
 v_c=1-\frac{Q_y'(1)}p+O(p^{-2}),\qquad
 \kappa_c=\ee^{-1}\left(1-\frac{Q_y(1)}p\right)+O(p^{-2}).
 \label{eq:fold-general}
\end{equation}
Thus, in the gamma case,
\begin{align}
 v_c&=1+\frac{\beta+1+\alpha}{p}+O(p^{-2}),\label{eq:fold-gamma-v}\\
 \kappa_c&=\ee^{-1}\left(1+\frac{\beta+(1+\alpha)/2}{p}\right)
                                          +O(p^{-2}).\label{eq:fold-gamma-k}
\end{align}
In the theta case,
\begin{align}
 v_c&=1+\frac{1-a_q(p)}p+O(p^{-2}),\label{eq:fold-q-v}\\
 \kappa_c&=\ee^{-1}\left(1+\frac{1/2-a_q(p)}p\right)
                                          +O(p^{-2}).\label{eq:fold-q-k}
\end{align}
The errors are uniform in the theta phase for fixed \(q\). Locally on
the two inverse branches,
\begin{equation}
 v=v_c\ \pm\sqrt{
      \frac{2(\kappa_c-\kappa)}{-\Psi_y''(v_c)}}
      +O(\kappa_c-\kappa),
 \qquad -\Psi_y''(v_c)\longrightarrow\ee^{-1}.
 \label{eq:sqrt}
\end{equation}
\end{theorem}
\begin{proof}
Differentiate \eqref{eq:psicorrection} near \(v=1\). Writing
\(v_c=1+d/p+O(p^{-2})\), the coefficient at order \(p^{-1}\) in
\(\Psi_y'(v_c)=0\) is \(-\ee^{-1}(d+Q_y'(1))\). This proves the first
part of \eqref{eq:fold-general}. The derivative of \(v\ee^{-v}\)
vanishes at one, so the displacement of the critical point contributes
only at order \(p^{-2}\) to the critical value. Evaluation of
\eqref{eq:psicorrection} at one proves the second part. Substitute
\eqref{eq:Q} for the displayed specializations. Finally, invert the
quadratic Taylor term at the nondegenerate critical point to obtain
\eqref{eq:sqrt}.
\end{proof}

The periodic coefficient is visible in both the inverse-gap law
\eqref{eq:q-gap} and the singularity law \eqref{eq:fold-q-k}.
Thus nonlinear inversion does not wash out the phase of the flat
ambiguity. The leading normal form is universal, but its first
correction remembers the realization scale.

\begin{corollary}[An abstract local criterion for the same radius law]
\label{cor:abstract}
The leading radius and Cayley-coefficient conclusions do not require a
moment model. Suppose a family of local equations has the form
\(H_y(t)=\delta yD_y(t)\), with \(H_y(t_0)=0\), \(D_y(t_0)>0\), and a
scale \(p_y\to\infty\). Suppose its data are real on the positive axis
and holomorphic on the expanding scaled discs, and, locally uniformly
in complex \(v\),
\begin{equation}
 \frac{p_yH_y(t_0(1+v/p_y))}{t_0}\longrightarrow v,
 \qquad
 \frac{D_y(t_0(1+v/p_y))}{D_y(t_0)}\longrightarrow\ee^v.
 \label{eq:abstract-fold}
\end{equation}
Then the branch through \(t_0\) has parameter radius
\(R_y\sim t_0/(\ee p_y yD_y(t_0))\), a real fold on its continuation,
and normalized fixed-order coefficient limits \(k^{k-1}/k!\).
If the convergence in \eqref{eq:abstract-fold} has rate \(O(p_y^{-1})\),
the coefficient errors have that rate as well.
\end{corollary}
\begin{proof}
The scaled deformation map is the quotient of the two expressions in
\eqref{eq:abstract-fold}. It converges holomorphically to \(v\ee^{-v}\).
The lower-radius continuation, local real critical point, upper-radius
obstruction, and coefficient extraction in the proof of
Theorem~\ref{thm:radius} apply without change.
\end{proof}

\section{Reproducible computations and verification boundaries}
\label{sec:computations}

The accompanying program \texttt{verify.py} uses exact integer
arithmetic for the cubic inverse coefficients and arbitrary-precision
arithmetic for finite analytic consistency tests. The recorded run used
Python 3.13.5, mpmath 1.3.0, and 130 decimal digits. All assertions in
that run passed. The exact environment and results are recorded in
\texttt{verification\_results.json}; there are no network dependencies.

\subsection{Exact arithmetic and independent transform evaluations}
The program computes the first twelve coefficients for \(m_n=(3n)!\)
using the finite coefficient formula and independently substitutes them
into the defining formal equation. The first five are displayed in
\eqref{eq:cubic-inverse}. For the theta family at \(q=1/2\), moments
through order twelve are checked separately on the even and odd
lattices. Transform defects are evaluated both as direct differences of
Stieltjes sums and from the theta formulas. The exact \(q^2\) equation
and the defect scaling are tested independently. The shifted-lattice
first-order equation and its alternating difference are also tested.

For the gamma defect, direct quadrature of the signed integral in the
cubic case gives the values in Table~\ref{tab:gamma}. The test compares
these with the negative of the closed form in \eqref{eq:cubic}.

\begin{table}[htbp]
\centering
\caption{Cubic signed Stieltjes defect, checked against the exact formula.}
\label{tab:gamma}
\begin{tabular}{@{}rl@{}}
\toprule
\(t\) & \(\displaystyle\int_0^\infty
 \frac{\ee^{-u}\sin(\sqrt3u-\pi/3)}{1+tu^3}\dd u=-D(t)\)\\
\midrule
\(1\)    & \(-0.141722777195878776146374\)\\
\(0.1\)  & \(-0.0303419687206032999616225\)\\
\(0.01\) & \(-0.000451920923529153552470793\)\\
\bottomrule
\end{tabular}
\end{table}

\subsection{Inverse transport and the fold}
Table~\ref{tab:qgap} compares the actual endpoint inverse gap with its
leading scale and first relative correction at \(q=1/2\). At integer
phases \(x=N\), symmetry gives \(P_h'(N-1/2)=0\), hence
\(a_q(N)=-1/2\). The final column is therefore
\(1-2y(N+1/2)\). The value at \(N=4\) is not in a particularly small
asymptotic regime; it is included to avoid displaying only favorable
large-\(N\) cases.

\begin{table}[htbp]
\centering
\caption{Theta-family inverse gap, \(q=1/2\) and \(y=2^{-N}\).}
\label{tab:qgap}
\begin{tabular}{@{}rlll@{}}
\toprule
\(N\) & \(D(y)\) & \((T_+-T_-)/(yD(y))\) & First correction\\
\midrule
4  & \(6.98194\cdot10^{-5}\)  & 0.644450288768 & 0.437500000000\\
8  & \(1.66462\cdot10^{-11}\) & 0.936987457980 & 0.933593750000\\
12 & \(6.05587\cdot10^{-23}\) & 0.993922761611 & 0.993896484375\\
20 & \(2.84745\cdot10^{-60}\) & 0.999960900309 & 0.999960899353\\
\bottomrule
\end{tabular}
\end{table}

At these integer phases,
\(D(y)/E_*(y)=0.004468440172262834\ldots\).
The phase is fixed, so this repeated ratio illustrates the periodic
formula rather than a general phase-independent constant.
For the cubic model, the measured ratio at \(y=10^{-4}\) is
\(0.991133461317105922\ldots\), compared with the first corrected value
\(0.990982261239872465\ldots\) from \eqref{eq:cubic-gap}.

To test the fold, the program fixes \(t_0=2^{-p}\) and computes
\(y=t_0/S_-(t_0)\). It evaluates the residual by subtracting the
integrands before summation, rather than subtracting two nearly equal
order-one numbers. The critical equation is
\begin{equation}
 H'(t)D(t)-H(t)D'(t)=0.
 \label{eq:numeric-fold}
\end{equation}
Table~\ref{tab:fold} compares the resulting critical values with
\eqref{eq:fold-q-k}. At these phases the predictor is
\(\ee^{-1}(1+1/p)\). The bounded last column is consistent with the
proved \(O(p^{-2})\) remainder.

\begin{table}[htbp]
\centering
\caption{The real theta fold and its first correction.}
\label{tab:fold}
\begin{tabular}{@{}rllll@{}}
\toprule
\(p\) & \(v_c\) & \(\kappa_c\) & \(\ee^{-1}(1+1/p)\)
& \(p^2\) times error\\
\midrule
12 & 1.15750238455 & 0.404091437651 & 0.398536061269 & 0.7999741990\\
20 & 1.08555562628 & 0.388065092191 & 0.386273413230 & 0.7166715845\\
32 & 1.05078622850 & 0.380048631028 & 0.379375673708 & 0.6891082956\\
\bottomrule
\end{tabular}
\end{table}

Finally, the exact first two parameter sectors at \(q=1/2,y=2^{-8}\)
are approximately
\[
 A_1=6.09270122025727\cdot10^{-14},\qquad
 A_2=7.19277265023981\cdot10^{-24}.
\]
For the endpoint displacement \(\delta=1\), the observed omitted tail
is approximately \(1.20633833887303\cdot10^{-33}\); the explicit bound
\eqref{eq:sector-tail}, evaluated with the theta disc, is approximately
\(1.42599129686661\cdot10^{-30}\). This verifies that the implemented
normalizations and the stated analytic estimate are mutually consistent
in this test case.

\subsection{What these computations certify}
The integer coefficient identities are exact finite calculations.
All quadratures, infinite-sum truncations, root computations, and
theta evaluations in the program use high-precision floating-point
arithmetic, not outward-rounded interval arithmetic. Thus the displayed
decimals are consistency evidence, not independently certified numerical
endpoints. The program does not prove the infinite identities, the
asymptotic theorems, or the absence of additional complex singularities.
Those assertions rely on the proofs above. A future interval
implementation would need explicit enclosures for every tail,
quadrature, and rounded arithmetic operation.

\section{Consequences for a transseries realization architecture}
\label{sec:architecture}

\subsection{Three objects that must not be identified}
The examples distinguish three different kinds of data. The first is a
formal coefficient sequence, or formal inverse \(\widehat T\). The
second is a selected analytic representative \(T_{a_0}\). The third is
the class of all admissible representatives and its flat deformations.
A symbolic inversion algorithm acts on the first object. An evaluated
residual certificate acts on the second. A theorem such as
\eqref{eq:information-floor} or \eqref{eq:radius} describes the third.
Confusing these objects turns a correct formal computation into an
incorrect claim of analytic uniqueness.

Positive measures provide a particularly useful test case because they
supply complete monotonicity, real-root uniqueness, and signed remainder
bounds, yet still leave nontrivial analytic moduli above the determinacy
threshold. The models therefore rule out the proposed implication
``positive Stieltjes structure plus all inverse coefficients implies a
unique analytic inverse,'' while identifying exactly when the
implication becomes valid.

\subsection{A natural representation of the additional data}
In the gamma example, a representation can retain an actual baseline
and a coefficient multiplying the generator
\(t^{-\beta}\ee^{-At^{-\alpha}}\). Nonlinear inversion generates
successive powers of that generator and its logarithmic derivatives.
The convergent parameter series gives a direct analytic meaning to
these extra sectors without assigning a sum to the divergent common
power series.

For theta, retaining only a constant coefficient at the log-Gaussian
scale loses information. A representation must also retain the phase
\(x\bmod1\), or the periodic function \(1/P_h(x-1/2)\). Its Fourier
expansion produces complex oscillatory factors, so a system restricted
to purely real nonoscillatory monomials should not silently treat it as
an ordinary one-term asymptotic equivalent. The fold correction shows
that this is not a superficial formatting issue: the phase affects the
location of an actual analytic singularity.

These observations are compatible with the repository's distinction
between formal well-based supports, flat-function interfaces, and
concrete analytic realization \cite{repo-series,repo-action}. They are
not a replacement for proving the missing general realization theorems.

\subsection{A realistic formalization boundary}
A staged formalization would begin with finite algebra and order facts:
the weak-composition coefficient identity, its sign, the factorial bound,
and the monotonicity/residual certificate. The moment-determinacy bridge
then needs measure push-forwards, polynomial approximation on a compact
interval, and uniqueness of Fourier transforms under exponential
integrability. The gamma and theta examples require complex integration
or equivalent special-function identities. The radius theorem requires
uniform holomorphic implicit inversion and a carefully stated local
fold lemma.

Existing repository modules on flatness, well-based supports, and
asymptotic scales are relevant interfaces, but their presence alone does
not certify these arguments. This report supplies no Lean source and
makes no claim that those dependencies have already been formalized in
the required form. The finite Python checks likewise are not a substitute
for a proof assistant.
% ed. (2026-09-29): the modules alluded to are named, with two related machine-checked facts (dossier of batch 51).
\begin{ednote}
The modules meant are
\path{Analysis/Transseries/Lean/Transseries/TransseriesFlat.lean},
\path{TransseriesWellBased.lean} and \path{TransseriesScale.lean} in the
same directory. Two machine-checked facts are related but not relied on:
\texttt{Fabius.exists\_eq\_in\_residual\_interval}
(\path{Analysis/FabiusFunction/Lean/FabiusFunction/MeanValueBracket.lean}),
the monotone residual bracket behind \eqref{eq:residual} in
Proposition~\ref{prop:inverse}, and \texttt{Fabius.cayleyTree}
(\path{Analysis/FabiusFunction/Lean/FabiusFunction/CayleyTreeFunction.lean}),
the real function \(-W_0(-w)\) for \(w\le\ee^{-1}\), without its power
series. No statement of this article is formalized.
\end{ednote}

\section{Further research questions}
\label{sec:questions}

\paragraph{Q1. Determine the full inverse ambiguity envelope.}
For a fixed indeterminate Stieltjes moment sequence, determine
\(\inf_\mu T_\mu(y)\) and \(\sup_\mu T_\mu(y)\) over all representing
probabilities, not only over the constructed affine families. Translate
the full moment parametrization into sharp nonlinear envelopes. Are the
even and odd lattice inverses asymptotically extremal at the least-bound
scale, or is the constant in the information floor strictly improvable?
No extremality claim for those two measures has been proved here.

\paragraph{Q2. Locate the exact nearest deformation singularities.}
Theorem~\ref{thm:radius} gives the leading radius, and
Theorem~\ref{thm:fold-correction} locates the real fold to one further
order. Determine whether that fold is the nearest singularity for all
sufficiently small \(y\). If not, classify the competing complex
singularities and their phase dependence. This is necessary before
identifying the second-order expansion of the Taylor radius with the
second-order expansion of the real fold.

\paragraph{Q3. Develop joint critical limits.}
The results keep \(\alpha,\beta,q\) fixed. Study the joint limit
\(\alpha\uparrow1/2\), \(y\downarrow0\), where
\(A=\sec(\pi\alpha)\) diverges and moment determinacy reappears at the
endpoint. Separately, analyze \(q\uparrow1\), \(y\downarrow0\), where
the theta oscillation is exponentially small in \(1/h\). A uniform
answer must distinguish small amplitude from exact absence of a phase.

\paragraph{Q4. Relate positive realizations to canonical summation.}
For the exact \(q\)-Euler family in
Proposition~\ref{prop:shifted}, identify which measures, if any, coincide
with particular analytic \(q\)-summation prescriptions. Determine the
additional sectorial growth or boundary data that restore uniqueness,
and then prove which of those selections commute with the nonlinear
map \(S\mapsto T\). The different summation procedures discussed in
\cite{divizio-zhang} make this a question about specified hypotheses,
not merely nomenclature.

\paragraph{Q5. Reconstruct the hidden parameter from minimal extra data.}
For the ordered one-parameter families, one exact transform value fixes
the parameter. Quantify the conditioning of this reconstruction when
that value is known only approximately. Combine the information floor
with the exact derivative \eqref{eq:parameter-derivative} to identify
optimal measurement scales. For the full moment-indeterminate family,
determine which finite or countable supplementary observations identify
a measure and its inverse.

\paragraph{Q6. Extend the radius criterion beyond the two models.}
Find verifiable conditions on logarithmic derivatives of a flat defect
that imply the complex-uniform limits in
Corollary~\ref{cor:abstract}. For example, formulate conditions for
regularly varying actions and for several competing flat scales.
Determine when one still obtains a single Lambert fold and when several
scales interact to produce a different limiting implicit equation.

\paragraph{Q7. Build an algebra of periodic-coefficient sectors.}
Construct a precise analytic or formal framework in which the
log-periodic coefficient \(1/P_h\), its derivatives, products, and
inverse-sector expansions are legitimate objects. Specify the support
and convergence conditions required for composition. Clarify the
boundary between a usable oscillatory calculus and a realization in an
ordered Hardy field; Proposition~\ref{prop:shifted} provides an explicit
test that any simultaneous-realization claim must pass.

\paragraph{Q8. Formalize the positive-measure bridge and certify numerics.}
A concrete first milestone is a formal proof of
Lemma~\ref{lem:coefficients} and the implication
``formal inverse coefficient bound \(\Rightarrow\) moment bound.''
A second is the exact identifiability theorem. Independently, implement
outward-rounded versions of the theta and gamma evaluations and use
\eqref{eq:residual} and \eqref{eq:sector-tail} to produce fully rigorous
inverse and sector enclosures. These are separate tasks from a full
formalization of the contour and fold arguments.

\section{Conclusion}
The positive Stieltjes inverse problem gives a complete, sharply scoped
answer to a realization question relevant to the ProveIt transseries
program. Formal inversion neither creates nor removes the
identifiability obstruction: that obstruction is exactly the underlying
moment problem. A sign-definite coefficient formula transfers the
classical order-two determinacy threshold to the nonlinear inverse.
Beyond it, exact exponential and theta defects expose the missing
analytic data.

The ambiguity is not an arbitrary nuisance term. It has an exact
parameter derivative, a convergent nonlinear sector expansion, an
information-theoretic error floor, and a universal rescaled deformation
map. The leading map \(v\ee^{-v}\) determines the asymptotic parameter
radius and Cayley coefficients, while the first fold correction retains
the specific gamma or theta structure. These conclusions supply
explicit theorems and counterexamples on which a broader realization
and certified-inversion theory can build.

\appendix
\section{The Lagrange identity used in both coefficient calculations}
\label{app:lagrange}
We record a short proof to make clear that the two uses of Lagrange
inversion in the paper have the same normalization. Let \(S(t)\) be a
formal series with \(S(0)=1\), and let \(T=yS(T)\). Formal substitution
\(y=t/S(t)\) in the residue formula for \([y^n]T\) gives
\begin{align*}
 [y^n]T
 &=\operatorname{res}_{t=0}
 \left(t^{-n}S(t)^n-t^{-(n-1)}S(t)^{n-1}S'(t)\right)\dd t\\
 &=[t^{n-1}]S(t)^n
   -\frac1n[t^{n-2}]\frac{\mathrm d}{\mathrm dt}S(t)^n\\
 &=\frac1n[t^{n-1}]S(t)^n.
\end{align*}
For \(n=1\), the coefficient with negative index is zero. For all
\(n\), the second term uses
\([t^{n-2}](S^n)'=(n-1)[t^{n-1}]S^n\).
The change-of-variables rule is valid for a formal substitution with
nonzero linear coefficient; it can also be checked directly on Laurent
monomials. In the analytic setting the same calculation follows from
Cauchy's integral formula on sufficiently small circles. Replacing
\(t\) by \(t_0+w\), and \(S\) by \(\phi_y\), yields the coefficient
formula of Theorem~\ref{thm:sectors}; a nonunit constant is allowed when
one uses the same residue argument with \(w=\delta\phi_y(w+t_0)\).

\section{Repository provenance and limits of the comparison}
\label{app:provenance}
The repository was inspected on 29 September 2026. The relevant live
path was \texttt{Analysis/Transseries}, not the older path under
\texttt{Analysis/FabiusFunction} used in some earlier research notes.
The retrieved Transseries subtree listing had Git tree identifier
\texttt{81058854a185ab903d9012f6de6f93da345b31d6}.
This identifies the inspected listing, not a claim that every file was
read or that the whole repository was frozen at that state.

The principal content sources were the transseries group README and
the README of the \emph{Action Accumulation and Nonlinear Inversion}
package \cite{repo-series,repo-action}. The former describes formal,
asymptotic, and analytic interfaces and their qualifications. The latter
provides the closest inspected comparison for hidden flat oscillations
and nonlinear transport, including an explicit warning about the general
Hahn-realization question. The canonical multi-volume sources and the
new research packages were not audited theorem by theorem.

An indexed repository code search for ``Stieltjes'' returned relevant
uses in other material; an indexed search for ``Wigert'' in the
Transseries path returned no matches. Neither search proves absence of
an equivalent theorem or historical novelty. The local
\texttt{SOURCE\_NOTES.md} records the inspected paths and the primary
literature used. No repository file was changed, and no external article is bundled in
this package.

\clearpage
\begin{thebibliography}{99}
\bibitem{repo-main}
V. Reshetnikov, \emph{ProveIt}, repository, inspected 29 September 2026.
\href{https://github.com/VladimirReshetnikov/ProveIt}{GitHub: VladimirReshetnikov/ProveIt}.

\bibitem{repo-series}
V. Reshetnikov, ProveIt contributors and research artifacts,
\emph{Series and transseries}, group README, inspected 29 September 2026.
\href{https://github.com/VladimirReshetnikov/ProveIt/blob/main/Analysis/Transseries/docs/series-and-transseries/README.md}{GitHub: Series and transseries README}.

\bibitem{repo-action}
ProveIt research artifact,
\emph{Action Accumulation and Nonlinear Inversion: A Laplace--measure
calculus, hidden oscillations, and limits of Hardy-field realization},
package README and editorial amendments, 29 September 2026.
\href{https://github.com/VladimirReshetnikov/ProveIt/blob/main/Analysis/Transseries/docs/series-and-transseries/Action_Accumulation_Nonlinear_Inversion/README.md}{GitHub: Action Accumulation and Nonlinear Inversion README}.

\bibitem{lin-stoyanov}
G. D. Lin and J. Stoyanov,
\emph{Moment determinacy of powers and products of nonnegative random
variables}, Journal of Theoretical Probability \textbf{28} (2015),
1337--1353.
\href{https://doi.org/10.1007/s10959-014-0546-z}{doi:10.1007/s10959-014-0546-z}.
Author manuscript: \href{https://arxiv.org/abs/1403.0301}{arXiv:1403.0301}.

\bibitem{chihara}
T. S. Chihara, \emph{A characterization and a class of distribution
functions for the Stieltjes--Wigert polynomials}, Canadian Mathematical
Bulletin \textbf{13} (1970), 529--532.
\href{https://doi.org/10.4153/CMB-1970-098-7}{doi:10.4153/CMB-1970-098-7}.

\bibitem{christiansen}
J. S. Christiansen, \emph{The moment problem associated with the
Stieltjes--Wigert polynomials}, Journal of Mathematical Analysis and
Applications \textbf{277} (2003), 218--245.
\href{https://doi.org/10.1016/S0022-247X(02)00534-6}{doi:10.1016/S0022-247X(02)00534-6}.

\bibitem{christiansen-koelink}
J. S. Christiansen and E. Koelink,
\emph{Self-adjoint difference operators and classical solutions to the
Stieltjes--Wigert moment problem}, Journal of Approximation Theory
\textbf{140} (2006), 1--26.
\href{https://doi.org/10.1016/j.jat.2005.11.010}{doi:10.1016/j.jat.2005.11.010}.

\bibitem{divizio-zhang}
L. Di Vizio and C. Zhang, \emph{On \(q\)-summation and confluence},
Annales de l'Institut Fourier \textbf{59} (2009), 347--392.
\href{https://doi.org/10.5802/aif.2433}{doi:10.5802/aif.2433}.
\href{https://arxiv.org/abs/0709.1610}{arXiv:0709.1610}.

\bibitem{nikolaev}
N. Nikolaev, \emph{Gevrey asymptotic implicit function theorem},
L'Enseignement Math\'ematique \textbf{70} (2024), 251--282.
\href{https://doi.org/10.4171/LEM/1061}{doi:10.4171/LEM/1061}.

\bibitem{dlmf}
F. W. J. Olver et al., eds., \emph{NIST Digital Library of Mathematical
Functions}, Chapter 17, especially Section 17.8, Jacobi's triple product.
\url{https://dlmf.nist.gov/17.8}. Consulted 29 September 2026.

% ed. (2026-09-29): the three entries below are editorial additions for the editorial note in Section 1.1.
\bibitem{ed:tai}
[Editorial addition, ProveIt, 2026-09-29.]
\emph{Transseries: the polynomial-logarithmic calculus, series reversal at
infinity, and inversion}, the canonical volume of the series-and-transseries
group:
\path{Analysis/Transseries/docs/series-and-transseries/Transseries_And_Inversion/transseries_and_inversion.tex}.

\bibitem{ed:ciq}
[Editorial addition, ProveIt, 2026-09-29.]
\emph{Certified Inversion Through the \(q\)-to-1 Transition}, research
package filed in batch~46:
\path{Analysis/Transseries/docs/series-and-transseries/Certified_Inversion_q_to_1_Transition/q_transition.tex}.

\bibitem{ed:ids}
[Editorial addition, ProveIt, 2026-09-29.]
\emph{A Spectral Representation for the Inverse Digamma Function}, research
package filed in batch~46:
\path{Analysis/Transseries/docs/series-and-transseries/Inverse_Digamma_Spectral_Representation/inverse_digamma_spectral.tex}.
\end{thebibliography}
\end{document}
